% CLAUDE REPORT::QCUBE-WR-0001::2026-11-15T17:06Z
% !TEX program = xelatex
% Compiles with XeLaTeX or pdfLaTeX (ASCII source).
\documentclass[11pt]{article}
\usepackage[margin=1in]{geometry}
\usepackage{amsmath,amssymb,amsthm}
\usepackage{xurl}
\newtheorem{theorem}{Theorem}
\newtheorem{lemma}[theorem]{Lemma}
\newtheorem{corollary}[theorem]{Corollary}
\theoremstyle{remark}
\newtheorem*{remark}{Remark}
\newcommand{\status}[1]{\textbf{[#1]}}
\newcommand{\Wb}{\overline W}

\title{CLAUDE QCUBE-WR-0001\\[2pt]
\large Worst-case minimum derivation cost on the hypercube-edge signed hypergraph:\\
an exact family where free optimal sequencing gives a linear $\Wb_R$, and a screen it implies}
\author{CLAUDE (scheduled routine; worst-case derivation directive)}
\date{2026-11-15 09:06 America/Los\_Angeles}

\begin{document}
\maketitle

\noindent\textbf{Status.} P versus NP remains \textbf{OPEN}. This is an additive report and not an LSPV
or TMR run. It changes no schedule, role, allocation or ownership region, and it leaves root
\texttt{main.tex} untouched. Every result is labelled with one of the project's strict statuses.

\section{Target and what this run attacks}
The shared directive asks for
\[
\Wb_R(N)=\max_{F\in\mathrm{UNSAT},\ |\langle F\rangle|\le N}\ \min_{\pi:F\vdash_R\bot}\mathrm{cost}_R(\pi),
\]
together with the weakest bridge $\mathrm{SAT}\in\mathrm P\Rightarrow\Wb_R(N)\le\mathrm{poly}(N)$.
The minimum ranges over all derivations, so the best sequence is granted for free and its discovery
is never charged. That is a legitimate lower-bound strategy.

This run takes up the directive's hypergraph formulation literally. It builds one explicit signed
hypergraph family on which the min--max can be computed exactly, and tests whether it can supply
hard instances (directive item 4). It cannot: free optimal sequencing makes the worst case linear.
That gives a quantitative screen for every future candidate family.

\section{CORPUS REVIEWED}
Read for this run: \texttt{AGENTS.md} (all headings, including the 2026-11-09 knot/link handoff
pointer and the claude world-ledger note); the 2026-11-14 fusion-Lipschitz five-agent refocus; the
TMR-1 current handoff (/047, tally/$\mathrm E\ne\mathrm{NE}$ family barrier); the root
\texttt{main.tex} header and the executive blocks of REFEREE296/297; REFEREE344 (balanced
p-optimal collapse); the LSPV-1/350 MUS worst-case quotient handoff; and the LSPV-1/398 Resolution
back-translation firewall. I also surveyed the existing claude reports (WORLD-LEDGER, TPORT,
ONETAPE) and ran repository-wide searches for \emph{p-optimal} (43 files), \emph{reflection
principle} (20), \emph{Pich} (159) and \emph{anticheck}.

\section{DUPLICATION / CONTRADICTION AUDIT}
\begin{itemize}
\item \emph{Bridge.} For a fixed Cook--Reckhow $R$, the bridge is equivalent to
$\mathrm P=\mathrm{NP}\Rightarrow R$ is p-bounded. For every $R$ that already has an unconditional
superpolynomial lower bound (Resolution by Haken; also Cutting Planes and bounded-depth Frege), the
bridge is therefore \emph{equivalent to} $\mathrm P\ne\mathrm{NP}$ and is not a weaker target. The
corpus already contains this in substance: REFEREE296 (relativization barrier and the conditional
size-optimal bridge), REFEREE317 and REFEREE344 (the p-optimal four-way equivalence), and
LSPV-1/398. I record it only as \status{KNOWN THEOREM / DUPLICATE} and claim no novelty.
\item \emph{MUS bounds.} LSPV-1/350 proved $\mu(F)-1\le W(F)\le 3^{\mu(F)-1}$ and an exact table for
$\le 3$ variables. This run reuses the lower bound $\mu-1$. Its new elements are the hypercube-edge
family, the linear upper bound $2\mu-1$ on that family, the fact that the bound $\mu-1$ is attained
only by tree-like proofs (Lemma~\ref{lem:tight}), and the exhaustive $n=4$ table. I found no
contradiction with the corpus.
\end{itemize}

\section{Exact objects}
\begin{theorem}[Hypercube-edge correspondence]\label{thm:corr}\status{PROVED/EXACT}
Fix $n\ge2$. A clause of width $n-1$ on $x_1,\dots,x_n$ omits exactly one variable $y$. It is
falsified by exactly the two assignments that differ only in $y$, which is an edge of $Q_n$. This map
is a bijection from width-$(n-1)$ clauses onto $E(Q_n)$; there are $n2^{n-1}$ of them, so the signed
hypergraph is the whole cube. A set $F$ of such clauses is UNSAT iff its edges cover $V(Q_n)$. It is
minimally UNSAT iff the cover is minimal, which holds iff every component is a star with at least two
vertices.
\end{theorem}
\begin{proof}
An assignment satisfies $F$ iff it falsifies no clause, that is, iff it is an uncovered vertex. A
cover is minimal iff every edge has an endpoint of degree~1. Since $Q_n$ is bipartite, this forces
every component to be a star. Conversely, a star forest that spans $V(Q_n)$ is a minimal cover.
\end{proof}

\begin{theorem}[Free optimal sequencing makes the family linear]\label{thm:lin}\status{PROVED/EXACT}
Let $F$ be UNSAT with $m$ clauses of width at least $n-c$. Then $m\ge 2^{n-c}$ and
$W_{\mathrm{Res}}(F)\le W_{\mathrm{TreeRes}}(F)\le 2^n-1\le 2^c m-1$. In particular, every MUS of
Theorem~\ref{thm:corr}, with $c=1$, satisfies
\[
 \mu-1\ \le\ W_{\mathrm{Res}}(F)\ \le\ W_{\mathrm{TreeRes}}(F)\ \le\ 2\mu-1 ,
\]
so on this family $W_{\mathrm{Res}}=\Theta(\mu)=\Theta(N/(\log N\log\log N))$, in particular $O(N)$, in
explicit encoded length.
\end{theorem}
\begin{proof}
Each clause falsifies at most $2^c$ of the $2^n$ points, and the clauses must jointly falsify every
point, so $m\ge 2^{n-c}$. The complete decision tree on all $n$ variables has $2^n$ leaves, and each
leaf, a single point, lies inside some falsified set. The standard conversion turns it into a
tree-like refutation with at most $2^n-1$ resolution steps. The lower bound $\mu-1$ is LSPV-1/350.
The encoded length is $N=\Theta(\mu\, n\log n)$, and $2^{n-1}\le\mu\le n2^{n-1}$ gives
$n=\Theta(\log N)$, hence $W=\Theta(\mu)=\Theta(N/(\log N\log\log N))$.
\end{proof}

\begin{corollary}[Volume screen for hard-instance candidates]\label{cor:screen}\status{PROVED/EXACT}
If a family $F_N$ has $W_{\mathrm{TreeRes}}(F_N)\ge N^{k}$ for every $k$, then
$n-w_{\min}(F_N)=\omega(\log N)$. A family whose clauses each falsify a $2^{-O(\log N)}$-dense part
of the cube is easy for tree-like Resolution, and so for every system that p-simulates it.
\end{corollary}
\begin{remark}
The corollary is a necessary screen, not a sufficient one. Proper 3-CNF with $n\to\infty$ always
passes it; Tseitin and the pigeonhole principle pass it, yet only Haken-type arguments give their
lower bounds. Families that the user's signed DAG hypergraph builds with near-full arity are retired
by it.
\end{remark}

\begin{lemma}[The bound $\mu-1$ is attained only by tree-like proofs]\label{lem:tight}\status{PROVED/EXACT}
Let $F$ be an MUS with $\mu$ clauses and let $\pi$ be a DAG Resolution refutation of $F$ with $I$
inferences. Then $I\ge\mu-1$, with equality only if $\pi$ is tree-like (every non-sink node has
out-degree~1). Hence $W_{\mathrm{Res}}(F)=\mu-1\iff W_{\mathrm{TreeRes}}(F)=\mu-1$. If
$W_{\mathrm{TreeRes}}(F)=\mu$, then $W_{\mathrm{Res}}(F)=\mu$.
\end{lemma}
\begin{proof}
Minimality forces all $\mu$ axioms to occur, because a refutation of a proper subset would contradict
minimality. Each inference has in-degree~2, and every node other than the sink has out-degree at
least~1. Counting edges gives $2I\ge \mu+I-1$, with equality iff every non-sink node has out-degree
exactly~1, which is the in-tree case. The last claim follows because $W_{\mathrm{Res}}\le
W_{\mathrm{TreeRes}}=\mu$ while $W_{\mathrm{Res}}\ne\mu-1$.
\end{proof}

\section{Computation}
\status{COMPUTATIONAL} All minimal edge covers were enumerated as spanning star forests and reduced
modulo the hyperoctahedral group (order $n!\,2^n$, i.e.\ 384 for $n=4$). Exact optimal tree-like
length was computed by dynamic programming over the $3^n$ subcubes. Each optimal tree was then
converted into an explicit weakening-free Resolution refutation and re-checked inference by inference
by a separate, self-contained checker, which also verifies UNSAT and clause-minimality.
\begin{center}\begin{tabular}{lrrl}
$n$ & MUS & classes & result\\\hline
2 & 2 & 1 & $W=\mu-1=1$\\
3 & 49 & 5 & $W_{\mathrm{Tree}}=W_{\mathrm{Res}}=\mu-1$ in all 5 classes ($\mu\in\{4,5,6\}$)\\
4 & 50{,}528 & 188 & see below; 4{,}143 checks, 0 failures, corrupted inference rejected
\end{tabular}\end{center}
For $n=4$, the 188 classes split as follows.
\begin{itemize}
\item 173 classes have $W_{\mathrm{Tree}}=\mu-1$, so $W_{\mathrm{Res}}=\mu-1$ exactly.
\item 13 classes have $W_{\mathrm{Tree}}=\mu$: eight with $\mu=9$, four with $\mu=10$ and one with
$\mu=11$. By Lemma~\ref{lem:tight}, $W_{\mathrm{Res}}=\mu$ exactly in all of them.
\item 2 classes, both perfect matchings with $\mu=8$, have $W_{\mathrm{Tree}}=9$ (orbit~32) and
$W_{\mathrm{Tree}}=11$ (orbit~8). Lemma~\ref{lem:tight} gives $W_{\mathrm{Res}}\in[8,9]$ and
$W_{\mathrm{Res}}\in[8,11]$. The exact DAG value is \textbf{OPEN}: the exhaustive backward DAG search
did not finish within its 30\,s chunk.
\end{itemize}
The worst case on $n=4$ is $\Wb_{\mathrm{Tree}}=11$, attained by the five $\mu=12$ classes, the
$\mu=11$ gap class and the $\mu=8$ matching. The matching is the only class whose gap above $\mu-1$
exceeds~2; it reaches~4. Its edges are $\{0,1\},\{2,6\},\{3,11\},\{4,12\},\{5,7\},\{8,10\},\{9,13\},\{14,15\}$,
two in each direction. These are isolated tree-versus-DAG test objects, not growing hard families,
since Theorem~\ref{thm:lin} caps the whole family linearly.

Reproduction (each invocation stays within the 30\,s rule; the residual DAG test is chunked per class):
\path{experiments/claude/qcube_tree_dp.py},
\path{experiments/claude/qcube_certify.py},
\path{experiments/claude/qcube_min_resolution.py},
\path{experiments/claude/qcube_dag_residual.py}. Raw data:
\path{results/raw/claude/2026-11-15_qcube_tree_n3.json},
\path{results/raw/claude/2026-11-15_qcube_tree_n4.json},
\path{results/raw/claude/2026-11-15_qcube_certify.jsonl},
\path{results/raw/claude/2026-11-15_qcube_dag_residual.jsonl}.

\section{Literature check (this run)}
\begin{itemize}
\item \emph{Res($\oplus$).} A superpolynomial size lower bound for unrestricted Res($\oplus$) is
still open as of 2026 (ECCC 2025/116; ECCC 2026/018). Depth-restricted bounds exist: superpolynomial
size for depth $o(n^2/\log^4 n)$ (ECCC 2026, Feb.) and the Itsykson--Knop supercritical size--depth
trade-off (ITCS 2026, LIPIcs vol.~362, art.~81). Consequence for directive item 4: a
hard instance certified only by Resolution does not yet survive ``algebraic/parity reasoning'' in
general. Only regular and depth-bounded parity reasoning is excluded today.
\item \emph{Bridge.} Pich--Santhanam, \emph{Towards P$\ne$NP from Extended Frege lower bounds}
(arXiv:2312.08163): P$\ne$NP follows if EF is not p-bounded \emph{and} one of these holds: feasible
SAT anticheckers, witnessing of NP$\not\subseteq$P/poly, or OWF from NP$\not\subseteq$P/poly plus
hardness of E. These are conditional bridges for $R=\mathrm{EF}$ in exactly the directive's
direction, and they are already cited 159 times in the corpus. Monroe, arXiv:2507.13576 v3, withdrew
its claim that no p-optimal proof system exists, so the existence of p-optimal systems stays
\textbf{OPEN}.
\item \emph{MUS structure.} Tarsi's lemma ($n\le\mu-1$) and the read-once refutability of
deficiency-1 MUS are consistent with the table: every $n=4$ class here has deficiency
$\mu-n\ge4$, and gaps appear only there.
\end{itemize}

\section{Barrier-transfer and post-quantum ledger}
\emph{Coding theory and covering.} Minimal edge covers of $Q_n$ are binary covering configurations of
radius~1 along one coordinate. The correspondence of Theorem~\ref{thm:corr} is an \textbf{EXACT
OBJECT TRANSFER}. No hardness transfers, because Theorem~\ref{thm:lin} is a volume argument.
\emph{VLSI area--time.} Volume and fan-out analogies are a \textbf{MISLEADING ANALOGY} here: the
obstruction is cube volume, not wire congestion. \emph{PQ item 4 (collision entropy).} Not
applicable: there is no proof-to-node map with nontrivial fibres in a tree-like optimum, so this is a
\textbf{FAILED TRANSFER}. \emph{Algorithmic exploitation.} The structure \emph{is} algorithmic: a
star-forest cover is decided and refuted in time linear in $2^n=O(N)$.

\section{WHAT CHANGED RELATIVE TO ROOT MAIN}
Root \texttt{main.tex} was not edited, because this routine owns no region in it. Proposed for
LSPV-5/TMR-5 adjudication: add Corollary~\ref{cor:screen} to the hostile-control ledger as a
necessary volume screen for any hard-instance family, and Lemma~\ref{lem:tight} as an exact
tree/DAG tightness criterion for MUS minima.

\section{ADOPTED / REJECTED / RETIRED PRIOR COMPONENTS}
Adopted: the free-optimal-sequence minimum, the LSPV-1/350 MUS quotient and its $\mu-1$ bound, and
strict status discipline. Retired as a hard-instance source \status{FALSIFIED/RETIRED}: signed
hypergraph families of near-full clause arity ($n-w_{\min}=O(\log N)$), including every
hypercube-edge family.

\section{STALE OR OVERSTATED CLAIMS FOUND}
None in the materials read. Calling the bridge ``weaker'' than P$\ne$NP is accurate only for proof
systems whose lower bounds are open, such as EF and Frege. For Resolution, Cutting Planes and
bounded-depth Frege the bridge is equivalent to P$\ne$NP.

\section{Open arrows}
A superpolynomial $\Wb_R$ for a fixed sufficiently strong $R$ is \textbf{OPEN}. The bridge
$\mathrm{SAT}\in\mathrm P\Rightarrow\Wb_R\le\mathrm{poly}$ for that same $R$ is \textbf{OPEN}. The
exact DAG values of the two $n=4$ perfect-matching classes are \textbf{OPEN} (computational).
P versus NP remains \textbf{OPEN}.

\end{document}
% END_OF_VERIFIED_REPORT::CLAUDE::QCUBE-WR-0001
